\chapter{Synthetic Noise Generation}
\label{ch:synthetic-noise}

\section{Purpose and Scope}
\label{sec:noise-purpose}

Trajectory simulations provide continuous, internally consistent descriptions of
flight, whereas a measurement system provides an incomplete and uncertain
record of that motion. Radar observations are affected by finite measurement
precision and by the availability and association of detections. A target may be
observed intermittently, and an accepted position measurement may differ from
the underlying position because of sensor and processing errors
\cite{barshalom1975clutter,barshalom2001estimation}. Additional errors can be
introduced as measurements are transmitted, stored, or processed. These
distinctions matter because missing measurements, ordinary measurement
uncertainty, and isolated data corruption do not affect a reconstructed
trajectory in the same way.

Previous classification research at Auburn University has examined the effect
of simulated radar measurement noise and filtering
\cite{eckert2020thesis}, as well as the sensitivity of classifications to
missing telemetry \cite{cervantes2024telemetry}. The present study adopts a
controlled \emph{synthetic degradation} approach. Clean numerical flight
histories are treated as reference trajectories, and selected portions of
their altitude histories are withheld or altered to produce the observations
available to subsequent analysis. The resulting measurement sets support
repeatable comparisons across degradation types and severity levels.

The objective is not to reproduce a specific radar installation in full.
Instead, the noise generator deliberately separates three mechanisms:
\emph{temporal sparsity}, \emph{additive altitude error}, and
\emph{digital-value corruption}. Sparsity is modeled differently during
powered and unpowered flight. Additive error perturbs the altitude values
that remain after sampling, and digital corruption introduces occasional
larger or irregular errors in those values. These mechanisms are applied
in sequence, so a degraded trajectory can contain both long gaps and
corrupted observations. The methodology below describes the physical or
data-handling motivation for each mechanism and the assumptions of its
synthetic implementation.

\section{Reference Trajectory and Observation Definition}
\label{sec:noise-reference}

Each test case begins with a simulated trajectory whose recorded time and
altitude are considered the reference solution. For consistency across
simulations, the reference trajectory is interpolated onto a uniform time
grid with a nominal interval of $\Delta t = 0.01~\mathrm{s}$. This grid is a
numerical basis for constructing comparable measurement sets; it is
\emph{not} a claim that a real radar necessarily reports measurements
at 100~Hz. Denoting the interpolated reference altitude by $h(t_i)$,
the grid is
\begin{equation}
    t_i=t_0+i\Delta t,\qquad i=0,\ldots,N-1.
    \label{eq:noise-time-grid}
\end{equation}

The reference history also provides the division between powered flight
and the subsequent coast phase. In the standard preparation procedure,
burnout is estimated from the simulated velocity history by identifying a
sustained reduction in acceleration, with a pre-apogee velocity-based
fallback when necessary. A known burnout time can alternatively be
supplied for an individual case. This reference information is used to
construct the phase-dependent \emph{observation schedule}; it is not
passed to the subsequent denoising procedure as known trajectory
information. In particular, the times at which observations are
retained are chosen by the synthetic experiment, not inferred from
a simulated radar-detection equation.

The current degradation model acts on \emph{altitude measurements}.
Time stamps remain exact, and the downrange history used in later
energy calculations is not degraded by this stage. Thus the
experiments measure sensitivity to altitude-data degradation,
not the complete uncertainty that would arise from errors in
all radar-derived coordinates and velocities. Reference altitude
values are retained separately for evaluation and are not supplied
as observations when the degraded trajectory is reconstructed.

\section{Temporal Sparsity and Interrupted Coverage}
\label{sec:noise-sparsity}

Missing observations arise whenever a tracking system fails to obtain,
associate, retain, or deliver a usable return. Low signal-to-noise
conditions, obscuration, tracking interruptions, and data-handling
limitations can contribute to incomplete histories. In clutter,
associating an observed return with the correct object is itself
an estimation problem \cite{barshalom1975clutter}. In the present
model, these possibilities motivate removing observations, but
their individual causes are not simulated separately.

Two sampling patterns are used because the \emph{temporal arrangement}
of the surviving observations matters in addition to their number.
A uniformly sparse history may still constrain the shape of the
whole trajectory. A history with the same number of observations
concentrated into short intervals can leave much longer periods
unconstrained. The model therefore treats powered and coast
coverage independently rather than deleting a single fixed
fraction of samples uniformly over the entire flight.

\subsection{Powered-Flight Sampling}
\label{sec:noise-powered}

During powered flight, a selected fraction of the reference altitude
samples is retained and distributed approximately uniformly in time.
Small random variations in the selected sample locations prevent
the simulated reports from falling on a perfectly regular
subsampling pattern. This construction represents an observing
system that maintains intermittent but relatively continuous
coverage while retaining only a limited set of reports.

Let $N_p$ be the number of points on the reference time grid during
powered flight and let $f_p$ be the powered-flight retention
fraction. The desired number of observations is approximately
\begin{equation}
    n_p \simeq \operatorname{round}(f_p N_p).
    \label{eq:noise-powered-count}
\end{equation}
The standard reference configuration uses $f_p=0.05$, so
approximately $5\%$ of powered-flight samples remain. Because
the selected measurements are deliberately spread across the
burn, this is \emph{not} an independent Bernoulli dropout
process: large gaps are less likely than they would be if
the same number of reports were selected entirely at random.

The choice to distribute observations across powered flight
is a modeling assumption made for controlled reconstruction
experiments. It is not a universal prediction about radar
coverage during propulsion.

\subsection{Coast-Flight Sampling}
\label{sec:noise-coast}

After the reference burnout time, retained observations are
organized into short contiguous \emph{bursts} separated by gaps.
The coast phase in this definition includes the remainder
of the trajectory, potentially including both ascent toward
apogee and subsequent descent. This pattern represents
interrupted tracking in which useful measurements arrive
during periods of contact but may be unavailable for
substantial intervals between contacts.

Let $N_c$ be the number of reference points after burnout,
with $f_c$ the coast retention fraction. The retained sample
count is approximately
\begin{equation}
    n_c \simeq \operatorname{round}(f_c N_c).
    \label{eq:noise-coast-count}
\end{equation}
Rather than distributing these $n_c$ points uniformly, the
generator assigns them to separated observation intervals.
The standard setup targets short bursts, with nominal
burst-duration controls of $0.3$--$0.8~\mathrm{s}$ and
a minimum target of five bursts when the retained sample
count permits. The actual grouping adapts to the requested
retention fraction and the number of available samples.
The intervals are placed at randomized locations without
overlap and with intervening missing-data regions.

The reference coast retention fraction is $f_c=0.01$,
corresponding to approximately $1\%$ of the potential
post-burnout samples. Even when the total retained fraction
is small, nearby samples within each burst contain useful
local trajectory information. By contrast, long gaps between
bursts require subsequent methods to infer motion over
intervals without observations. This difference is the
principal reason for using burst-like sampling in addition
to simple sparse sampling.

For the primary noise configuration, no observations are
guaranteed near apogee and no minimum-total-observation
backfill is imposed. Optional experimental safeguards can
retain observations within a specified apogee window or add
samples when a measurement set is too small, but those
safeguards are disabled in the standard robustness setting.
Consequently, severe data loss can remain severe rather
than being silently corrected by the generator.

\section{Additive Altitude Measurement Error}
\label{sec:noise-additive}

A retained measurement can be inaccurate even when it has
been detected and associated successfully. Radar-derived
position uncertainty may reflect receiver noise, propagation
effects, angular and range estimation errors, and subsequent
coordinate transformation or processing. Such effects are
frequently represented with statistical measurement-error
models \cite{barshalom2001estimation,eckert2020thesis}.
Here, they motivate adding a random error directly to
the altitude value of each \emph{retained} sample.

The model uses a zero-mean Laplace distribution:
\begin{equation}
    \widetilde{h}_i=h(t_i)+\epsilon_i,\qquad
    \epsilon_i\sim\operatorname{Laplace}(0,b),
    \label{eq:noise-additive-model}
\end{equation}
with probability density
\begin{equation}
    p(\epsilon)=\frac{1}{2b}
       \exp\left(-\frac{|\epsilon|}{b}\right),
    \qquad b>0.
    \label{eq:noise-laplace-pdf}
\end{equation}
The errors are drawn independently for the surviving
samples. The distribution is symmetric, so positive
and negative altitude errors are equally likely.
The parameter $b$ controls its scale, with
\begin{equation}
    \mathbb{E}[\epsilon]=0,\qquad
    \operatorname{SD}(\epsilon)=\sqrt{2}\,b.
    \label{eq:noise-laplace-spread}
\end{equation}

A Laplace distribution is used as a \emph{stress model}
because its heavier tails produce occasional larger
errors in addition to typical small-to-moderate
perturbations. This choice is useful when assessing
whether a reconstruction method depends excessively
on each individual observation. It should not be
interpreted as a finding that radar-altitude errors
are universally Laplace-distributed. Establishing
a sensor-specific error law would require the
characteristics or measurements of that sensor.

The nominal additive scale is $b=1000~\mathrm{m}$;
its corresponding standard deviation is approximately
$1414~\mathrm{m}$. All trajectory altitudes and error
amplitudes are expressed in meters. Although a
legacy software parameter name contains a feet
suffix, the current procedure performs no unit
conversion and adds this error directly to
meter-valued altitude data. This distinction is
essential for reproducibility of the reported
noise levels.

\section{Digital Corruption and Isolated Outliers}
\label{sec:noise-bit-errors}

A second type of value error is introduced \emph{after}
additive measurement noise. Unlike an ordinary sensor
measurement perturbation, this mechanism represents
possible errors introduced while numerical measurements
are encoded, transported, or stored. Bit corruption is
therefore a model of \emph{digital-data integrity},
not an assertion that radar propagation itself flips
bits in a floating-point altitude value.

Each retained noisy altitude is represented in
32-bit floating-point form. Such a value contains
a sign, an exponent, and a 23-bit stored fraction
(or mantissa). For a randomly selected fraction of
the retained samples, the model reverses between one
and three \emph{distinct fraction bits}. Sign and
exponent bits are left unchanged. The modified bit
pattern is then interpreted as an altitude value.

The standard corruption probability is $f_b=0.10$,
so approximately $10\%$ of \emph{surviving
observations}, not $10\%$ of every point on the
original reference grid, undergo this operation.
Because the bit positions are selected randomly,
the resulting numerical error is not of fixed
magnitude. Changing a low-significance fraction
bit may produce a small perturbation, while changing
a high-significance fraction bit may produce a
substantially larger error. The same number of
bit changes also has different absolute effects
on numbers with different exponents.

This mechanism consequently produces irregular
outliers that are distinct from the independently
drawn additive errors. It is intentionally
constrained to the fraction bits rather than the
sign or exponent fields: the latter could produce
more extreme or nonphysical numerical changes.
The model tests sensitivity to corrupted stored
values without claiming to reproduce the
error-correction behavior of a particular
telemetry or computing system.

\section{Sequence of Degradation and Reference Conditions}
\label{sec:noise-sequence}

All three mechanisms are applied to the same
reference altitude history in a fixed order:
\begin{equation}
\begin{gathered}
\text{Reference trajectory}
\ \longrightarrow\
\text{phase-dependent observation selection}
\\
\longrightarrow\
\text{additive altitude error}
\ \longrightarrow\
\text{selected floating-point bit corruption}
\\
\longrightarrow\
\text{degraded altitude observations}.
\end{gathered}
\label{eq:noise-workflow}
\end{equation}
Observation selection comes first: missing samples
are omitted, rather than assigned artificial
numerical values. The Laplace perturbations are
then added only to the observations that remain.
Digital corruption acts on a randomly chosen
subset of those noisy observations. Thus one
retained measurement may contain both additive
error and bit corruption, while an omitted
measurement is never artificially restored by
a later stage.

Table~\ref{tab:noise-reference-settings}
summarizes the nominal configuration used in
the energy-classification noise experiments.
Each parameter can be varied independently.
The two retention fractions govern
\emph{availability}; the Laplace scale governs
ordinary measurement perturbations and their
tails; and the bit-corruption fraction governs
the frequency of digital outliers.

\begin{table}[htbp]
    \centering
    \caption{Reference configuration of the synthetic altitude degradation model.}
    \label{tab:noise-reference-settings}
    \begin{tabular}{lll}
        \hline
        Quantity & Symbol & Reference setting \\
        \hline
        Reference grid interval & $\Delta t$ & $0.01~\mathrm{s}$ \\
        Powered-flight retention & $f_p$ & $0.05$ ($5\%$) \\
        Coast-flight retention & $f_c$ & $0.01$ ($1\%$) \\
        Laplace altitude-error scale & $b$ & $1000~\mathrm{m}$ \\
        Fraction of retained values corrupted & $f_b$ & $0.10$ ($10\%$) \\
        Distinct fraction bits per corrupted value & -- & 1--3 \\
        Nominal coast-burst durations & -- & $0.3$--$0.8~\mathrm{s}$ \\
        Protected apogee window & -- & Disabled \\
        Minimum-observation backfill & -- & Disabled \\
        \hline
    \end{tabular}
\end{table}

The random sampling locations, additive errors,
and bit selections are generated from controlled
random-number states. Repeating an experiment with
the same source trajectory, parameters, and random
seed reproduces the same degraded observations.
Different random seeds generate independent
realizations of the specified degradation condition.
This permits performance to be examined across
multiple possible gap and error arrangements,
rather than relying on a single fortunate
or unfavorable realization.

\section{Degradation Severity and Experimental Comparisons}
\label{sec:noise-severity}

The noise model is intended to support both
\emph{isolated-mechanism} and \emph{combined}
stress studies. Isolated studies change the
powered/coast retention settings, the additive
error scale, or the digital corruption rate
while keeping the other mechanisms at their
clean settings. In this way, changes in
downstream performance can be associated
more directly with a particular source
of degradation. The powered and coast
retention fractions are varied together
in the illustrated sparsity sweep, but
remain separate parameters of the model.

Combined conditions vary the three mechanisms
at once to examine their interaction. The
noise-study configuration employs the
conditions in Table~\ref{tab:noise-combined-levels}.
The labels describe progressively stronger
\emph{synthetic} stress conditions rather
than standardized categories of real radar
performance.

\begin{table}[htbp]
    \centering
    \caption{Combined synthetic-degradation conditions used in the noise study.}
    \label{tab:noise-combined-levels}
    \begin{tabular}{lrrrr}
        \hline
        Condition & $f_p$ & $f_c$ & $b$ (m) & $f_b$ \\
        \hline
        Clean      & 1.00 & 1.000 & 0    & 0.00 \\
        Light      & 0.50 & 0.250 & 250  & 0.01 \\
        Moderate   & 0.25 & 0.100 & 500  & 0.03 \\
        High       & 0.10 & 0.030 & 1000 & 0.05 \\
        Very high  & 0.05 & 0.010 & 1000 & 0.10 \\
        Extreme    & 0.02 & 0.005 & 2000 & 0.20 \\
        \hline
    \end{tabular}
\end{table}

The isolated additive-error study uses scales of
$0$, $100$, $250$, $500$, $1000$, and
$2000~\mathrm{m}$; the isolated digital-error
study considers corruption fractions of
$0$, $0.01$, $0.03$, $0.05$, $0.10$, and
$0.20$. The isolated sparsity conditions
range from full retention to $2\%$ powered
and $0.5\%$ coast retention. Multiple
random realizations at each condition
allow variation caused by the random
placement of gaps or corrupted values
to be distinguished from the effect of
the chosen severity level.

The clean condition is an essential
control: it retains all reference points
and applies neither additive noise
nor bit corruption. Comparisons with
this condition distinguish the effect
of data degradation from differences
already present between the underlying
simulated trajectories.

\section{Model Assumptions and Implications}
\label{sec:noise-limitations}

The synthetic construction makes several
deliberate simplifications. First, it
does not compute detection probability
from range, radar cross section,
antenna geometry, or signal-to-noise
ratio. The distinct powered and coast
coverage patterns are prescribed test
conditions, not consequences of an
explicit sensor model. Second, the
additive altitude errors are independent
and have zero mean: systematic
calibration offsets, temporally
correlated drift, and propagation biases
are not separately modeled. Third,
the bit-flip mechanism is a selected
digital-value corruption test and
does not incorporate error detection
or correction during transmission.
Finally, only altitude is corrupted,
so the study does not represent
the full propagation of coupled
position-coordinate errors into
all reconstructed motion quantities.

These choices limit claims about
operational realism, but also make
the degradation mechanisms
interpretable and repeatable.
They establish a controlled set of
observations for asking how much
trajectory information can be
recovered when the data are
temporally incomplete, imprecise,
and occasionally corrupted.
The reconstruction of continuous
histories from those observations,
and the handling of outliers during
that reconstruction, are considered
in the following chapter.
